\documentclass[11pt]{article}
\usepackage[a4paper,margin=1in]{geometry}
\usepackage{amsmath,amssymb}
\usepackage{array}
\usepackage{hyperref}

\newcommand{\dd}{\mathrm{d}}
\newcommand{\vt}{v_t}
\newcommand{\Kn}{\mathrm{Kn}}
\newcommand{\sgn}{\operatorname{sgn}}
\newcommand{\avg}[1]{\langle #1 \rangle}
\newcommand{\bu}{\mathbf{u}}
\newcommand{\bv}{\mathbf{v}}
\newcommand{\bq}{\mathbf{q}}
\newcommand{\bx}{\mathbf{x}}
\newcommand{\bE}{\mathbf{E}}
\newcommand{\bB}{\mathbf{B}}

\title{Closure models: a high-level tutorial\\[4pt]
\large From gases and turbulence to plasmas}
\author{Working note}
\date{16 November 2026}

\begin{document}
\maketitle

\noindent\emph{Status.} This tutorial is written from general knowledge of
the subject. The equations are standard, but the text and the bibliography
were not checked against the cited sources, with the exception of the
Hammett--Perkins material in Section~\ref{sec:collisionless}, which is
derived and checked in the companion note
\texttt{hammett\_perkins\_closure\_derivation.tex}.

\tableofcontents

\section{What a closure problem is}

Suppose a system has a full state $X$ obeying $\dot X = F(X)$, and we only
want to track a reduced state $x = \mathcal{P}X$: an average, a few moments,
a coarse-grained field. Applying $\mathcal{P}$ to the dynamics gives
\begin{equation}
  \dot x = \mathcal{P}F(X).
  \label{eq:unclosed}
\end{equation}
The right-hand side depends on the full state $X$, not only on $x$. That is
the closure problem. A \emph{closure} is a rule that replaces
$\mathcal{P}F(X)$ by something computable from $x$ alone,
\begin{equation}
  \dot x = f(x) + \underbrace{\bigl[\mathcal{P}F(X) - f(x)\bigr]}_{\text{closure defect}} ,
\end{equation}
and the term in brackets is what the closure gets wrong. Every scale
transition in this project can be described with five items: the exact
hierarchy, the limiting process, the retained state, the discarded
information, and the closure defect.

The problem always has one of two causes, and usually both:
\begin{itemize}
  \item \textbf{Nonlinearity.} The average of a product is not the product
    of averages.
  \item \textbf{Transport.} Streaming or advection couples each retained
    quantity to one that is not retained, producing an infinite hierarchy.
\end{itemize}

\section{Three places the problem first appeared}

\subsection{Kinetic theory of gases}

A dilute gas is described by a distribution $f(\bx,\bv,t)$ obeying the
Boltzmann equation,
\begin{equation}
  \partial_t f + \bv\cdot\nabla f = C[f],
\end{equation}
where $C$ is the collision operator. Take velocity moments: density
$n = \int f\,\dd^3 v$, flow $n\bu = \int \bv f\,\dd^3 v$, and temperature
$\tfrac{3}{2}nT = \int \tfrac{1}{2}m|\bv-\bu|^2 f\,\dd^3 v$ (with
$k_B = 1$). Because collisions conserve particles, momentum and energy, the
moment equations are
\begin{align}
  \partial_t n + \nabla\cdot(n\bu) &= 0, \\
  mn\,\frac{D\bu}{Dt} &= -\nabla p - \nabla\cdot\boldsymbol{\pi}, \\
  \frac{3}{2}\,n\,\frac{DT}{Dt} + p\,\nabla\cdot\bu
     &= -\nabla\cdot\bq - \boldsymbol{\pi}:\nabla\bu ,
\end{align}
with $p = nT$. These are exact, and they are not closed: the viscous stress
$\boldsymbol{\pi}$ and heat flux $\bq$ are higher moments of $f$. Writing
equations for them brings in still higher moments. The term
$\bv\cdot\nabla f$ is responsible: multiplying by $\bv^k$ and integrating
always produces a moment of order $k+1$. Maxwell wrote this hierarchy down in
1867; the systematic closures came from Hilbert, Chapman and Enskog
(1912--1917) and from Grad (1949).

\subsection{Turbulence}

For an incompressible fluid, split the velocity into a mean and a
fluctuation, $u_i = \bar u_i + u_i'$, and average the Navier--Stokes
equations (Reynolds, 1895):
\begin{equation}
  \partial_t \bar u_i + \bar u_j\,\partial_j \bar u_i
  = -\frac{1}{\rho}\,\partial_i \bar p + \nu\,\nabla^2 \bar u_i
    - \partial_j \avg{u_i' u_j'} .
\end{equation}
The Reynolds stress $\avg{u_i'u_j'}$ is not a function of the mean flow.
Its own evolution equation contains triple correlations, and so on. Here the
cause is the quadratic nonlinearity. Large-eddy simulation has the same
structure with a spatial filter in place of the average.

\subsection{Many-body statistical mechanics}

For $N$ interacting particles, the reduced $s$-particle distributions obey
the BBGKY hierarchy, schematically
\begin{equation}
  \partial_t f_s = \mathcal{L}_s f_s + \mathcal{C}_{s,s+1} f_{s+1},
  \qquad s = 1, 2, \dots, N-1,
\end{equation}
where $\mathcal{L}_s$ is the dynamics of $s$ particles among themselves and
$\mathcal{C}_{s,s+1}$ is their interaction with one more. A kinetic equation
is a closure of the $s=1$ equation: the Boltzmann equation follows from
assuming that particles about to collide are uncorrelated,
$f_2 \approx f_1 f_1$, in the low-density limit.

\subsection{The common pattern}

\begin{center}
\renewcommand{\arraystretch}{1.4}
\begin{tabular}{>{\raggedright}p{3.2cm}>{\raggedright}p{3.6cm}>{\raggedright}p{3.4cm}p{3.6cm}}
  Problem & Retained state & Unclosed quantity & Discarded information \\ \hline
  Kinetic to fluid & $n$, $\bu$, $T$ & $\boldsymbol{\pi}$, $\bq$
     & shape of $f$ in velocity \\
  Turbulence & $\bar{\bu}$, $\bar p$ & $\avg{u_i'u_j'}$
     & fluctuations below the filter \\
  Particles to kinetic & $f_1$ & $f_2$ & correlations between particles \\
\end{tabular}
\end{center}

\section{The standard ways to close}

\subsection{Asymptotic closure}

If a small parameter makes the discarded variables relax quickly to a state
determined by the retained ones, the closure can be derived. For a gas the
parameter is the Knudsen number,
\begin{equation}
  \Kn = \frac{\lambda}{L},
\end{equation}
the mean free path over the gradient length. The Chapman--Enskog expansion
writes $f = f_M\,(1 + \Kn\,\phi_1 + \Kn^2\phi_2 + \dots)$ about the local
Maxwellian $f_M(n,\bu,T)$ and solves for $\phi_1$. At zeroth order
$\boldsymbol{\pi} = \bq = 0$ (the Euler equations). At first order,
\begin{equation}
  \boldsymbol{\pi} = -\mu\left(\nabla\bu + \nabla\bu^{\mathsf T}
     - \tfrac{2}{3}(\nabla\cdot\bu)\,\mathbf{I}\right),
  \qquad
  \bq = -\kappa\,\nabla T ,
  \label{eq:NSF}
\end{equation}
which are the Navier--Stokes--Fourier equations, with
$\mu \sim mn\,\vt\lambda$ and $\kappa \sim n\,\vt\lambda$ computed from the
collision operator. The closure is local in space and time, and its error is
controlled by $\Kn$.

\subsection{Closure by an ansatz for the distribution}

Instead of expanding, assume a shape for $f$ with as many free parameters as
retained moments, then compute the missing moment from that shape.
\begin{itemize}
  \item \textbf{Grad (1949):} $f = f_M \times$ (a polynomial in $\bv$), with
    stress and heat flux promoted to dynamical variables (13 moments).
  \item \textbf{Maximum entropy (Levermore 1996):}
    $f = \exp\bigl(\sum_i \alpha_i m_i(\bv)\bigr)$, the least-biased
    distribution consistent with the retained moments.
\end{itemize}
These give evolution equations for the fluxes, so they have finite relaxation
times, but they can lose hyperbolicity or realizability away from
equilibrium.

\subsection{Gradient-diffusion models}

When no small parameter is available, postulate the form of \eqref{eq:NSF}
and model the coefficient. For turbulence (Boussinesq, 1877),
\begin{equation}
  -\avg{u_i'u_j'} + \tfrac{2}{3}k\,\delta_{ij}
  = \nu_t\left(\partial_i \bar u_j + \partial_j \bar u_i\right),
\end{equation}
with the eddy viscosity $\nu_t$ supplied by a mixing length
($\nu_t = \ell_m^2\,|\partial \bar u/\partial y|$, Prandtl 1925), by
transport equations for turbulence quantities ($k$--$\varepsilon$), or by the
filter width ($\nu_t = (C_s\Delta)^2|\bar S|$, Smagorinsky 1963). These are
models, not limits: turbulent eddies have no scale separation from the mean
flow that would justify them.

\subsection{Memory closures}

The Mori--Zwanzig formalism projects the full dynamics onto the retained
variables without approximation. For retained variables $a(t)$ the result is
a generalized Langevin equation,
\begin{equation}
  \frac{\dd a}{\dd t} = i\Omega\,a(t)
     - \int_0^t K(s)\,a(t-s)\,\dd s + F(t),
\end{equation}
with a Markovian term, a memory kernel $K$, and a term $F$ that depends on
the discarded initial data. This shows what any local closure gives up:
memory and noise. Asymptotic closures correspond to a kernel that decays much
faster than $a$ evolves.

\subsection{Learned closures}

Fit the unclosed term, or a correction to a classical closure, from
higher-fidelity data. The learned object can be a local function, a nonlocal
operator, or a memory kernel. The structural questions are the same as for
any other closure.

\subsection{How a closure is judged}

\begin{itemize}
  \item \textbf{Conservation} of mass, momentum and energy.
  \item \textbf{Realizability:} the closed moments correspond to some
    non-negative distribution.
  \item \textbf{Well-posedness and stability}, including an entropy
    inequality where one exists.
  \item \textbf{Limit recovery:} the closure reduces to the trusted result
    where the small parameter is small.
  \item \textbf{A posteriori behaviour:} accuracy when embedded in a solver
    over long times, not only accuracy of the closure term in isolation.
\end{itemize}

\section{Why neutral gases are the easy case}

Collisions between neutral molecules are short-range and frequent. Three
facts follow, and each will fail for plasmas.
\begin{enumerate}
  \item \textbf{Fast local relaxation.} Collisions drive $f$ to a local
    Maxwellian in a few collision times, so the kinetic state is determined
    by $n$, $\bu$, $T$ and their gradients.
  \item \textbf{One small parameter.} All molecules have comparable mean
    free paths, so a single $\Kn$ controls the expansion.
  \item \textbf{Locality.} Fluxes depend on gradients at the same point and
    time.
\end{enumerate}

\section{What makes plasmas different}

\subsection{Long-range forces and collective behaviour}

The Coulomb force is long-range, so each particle interacts with many others
simultaneously. The relevant parameter is the number of particles in a Debye
sphere, $\Lambda \sim n\lambda_D^3$. When $\Lambda \gg 1$, the leading-order
closure of BBGKY is not a collision term at all but a mean field:
\begin{equation}
  \partial_t f_s + \bv\cdot\nabla f_s
  + \frac{q_s}{m_s}\left(\bE + \bv\times\bB\right)\cdot\nabla_{\bv} f_s
  = C[f],
\end{equation}
with $\bE$ and $\bB$ determined self-consistently by Maxwell's equations.
Collisions enter as a correction of relative size $\sim \ln\Lambda/\Lambda$,
and they are the cumulative effect of many small-angle deflections, described
by a Fokker--Planck (Landau) operator, not a Boltzmann operator. The
expansion parameter for the particle-to-kinetic closure is weak coupling, not
low density.

When the potential energy between neighbours is comparable to or larger than
the thermal energy (coupling parameter $\Gamma \gtrsim 1$, as in dusty
plasmas), $\Lambda$ is not large, pair correlations are not small, and this
kinetic closure fails as well.

\subsection{Collisions weaken with energy}

The Coulomb collision frequency falls steeply with speed,
\begin{equation}
  \nu(v) \propto \frac{n\ln\Lambda}{v^3},
  \qquad
  \lambda(v) = \frac{v}{\nu(v)} \propto v^4 .
\end{equation}
There is no single mean free path. The Chapman--Enskog result for a plasma
(Spitzer--H\"arm 1953; Braginskii 1965) is $\bq = -\kappa\nabla T$ with
$\kappa \propto T^{5/2}$, but the heat is carried by electrons several times
faster than thermal, whose mean free path is far longer than the thermal one.
The local closure therefore fails when the thermal mean free path is still
only of order one percent of the temperature gradient length. Two standard
repairs are used in practice:
\begin{align}
  \text{flux limiter:}\quad
    & q = \min\bigl(q_{\mathrm{SH}},\; f\,n T \vt\bigr),
      \quad f \text{ tuned}, \\
  \text{nonlocal kernel:}\quad
    & q(x) = \int W(x,x')\,q_{\mathrm{SH}}(x')\,\dd x' ,
\end{align}
the second due to Luciani, Mora and Virmont (1983) and developed into
multigroup models such as Schurtz--Nicola\"i--Busquet (2000). The closure has
become a nonlocal functional of the temperature profile.

\subsection{No collisions does not mean no closure physics}
\label{sec:collisionless}

In hot, dilute plasmas the mean free path exceeds the system size. Nothing
then relaxes $f$ toward a Maxwellian, and Chapman--Enskog has no small
parameter. What takes the place of collisions is phase mixing: particles
streaming at different speeds smear out perturbations, and particles resonant
with a wave ($v = \omega/k$) exchange energy with it (Landau damping).

The Hammett--Perkins closure builds this into a fluid model. Compare the two
heat-flux closures in Fourier space:
\begin{equation}
  \text{collisional:}\quad q_k = -\kappa\,ik\,T_k,
  \qquad\qquad
  \text{collisionless:}\quad
  q_k = -n_0\sqrt{\frac{8}{\pi}}\;\vt\;i\,\sgn(k)\,T_k .
\end{equation}
The collisional closure damps temperature perturbations at a rate
$\propto k^2$; the collisionless one at a rate $\propto |k|\vt$, the
phase-mixing rate. In real space the first is a derivative and the second is
a Hilbert transform, a nonlocal operator with no intrinsic length scale. It
is matched to the linear kinetic response and is exact only in limiting
cases.

\subsection{The magnetic field makes the closure anisotropic}

A magnetic field confines particles to gyro-orbits of radius $\rho$ across
the field while leaving them free along it. This introduces a second small
parameter, $\rho/L$, that works even without collisions.
\begin{itemize}
  \item \textbf{Across the field,} gyromotion provides locality. With heat
    fluxes neglected, the pressure tensor is gyrotropic and obeys the
    double-adiabatic (Chew--Goldberger--Low 1956) laws
    \[
      \frac{\dd}{\dd t}\!\left(\frac{p_\perp}{nB}\right) = 0,
      \qquad
      \frac{\dd}{\dd t}\!\left(\frac{p_\parallel B^2}{n^3}\right) = 0 .
    \]
  \item \textbf{Along the field} there is no such help, and the parallel
    closure is where the nonlocal physics of the previous two subsections
    enters.
  \item \textbf{With collisions,} transport coefficients differ by many
    orders of magnitude between directions:
    $\kappa_\perp/\kappa_\parallel \sim (\Omega\tau)^{-2}$, where $\Omega$
    is the gyrofrequency and $\tau$ the collision time.
\end{itemize}

\subsection{More fields and more scales}

A plasma has at least two species with a mass ratio of 1836 or more, which
can have different temperatures and flows, and electromagnetic fields that
are dynamical variables in their own right. The Debye length, the skin
depths and the gyroradii of each species are all distinct scales, and a
closure valid at one may fail at another.

\subsection{Closures in sequence}

A plasma model typically stacks several closures, each with its own defect:
\begin{center}
\renewcommand{\arraystretch}{1.4}
\begin{tabular}{>{\raggedright}p{4.2cm}>{\raggedright}p{4.6cm}p{5.2cm}}
  Transition & Small parameter & Fails when \\ \hline
  Particles $\to$ kinetic & $1/\Lambda$ (weak coupling)
     & strong coupling, dense or dusty plasmas \\
  Kinetic $\to$ fluid, collisional & $\lambda/L$
     & steep gradients, fast particles \\
  Kinetic $\to$ fluid, magnetized & $\rho/L$
     & unmagnetized regions, parallel dynamics \\
  Fluid $\to$ mean transport & none in general
     & always a model (turbulence) \\
\end{tabular}
\end{center}

\section{Side-by-side summary}

\begin{center}
\renewcommand{\arraystretch}{1.5}
\begin{tabular}{>{\raggedright}p{2.9cm}>{\raggedright}p{3.4cm}>{\raggedright}p{3.4cm}p{4.3cm}}
  & Neutral gas & Fluid turbulence & Plasma \\ \hline
  What relaxes the discarded variables
    & frequent short-range collisions
    & nothing fast; cascade
    & weak, speed-dependent collisions; phase mixing; gyromotion \\
  Small parameter
    & $\Kn$
    & none
    & several ($1/\Lambda$, $\lambda/L$, $\rho/L$), often not small \\
  Form of closure
    & local gradient law
    & modelled eddy viscosity
    & nonlocal in space, anisotropic, sometimes with memory \\
  Status
    & derived, with error control
    & modelled
    & derived only in limits; modelled or matched elsewhere \\
\end{tabular}
\end{center}

\paragraph{In one sentence.} In a gas the closure is justified because
collisions enforce local equilibrium below every scale that is resolved; in a
plasma that justification is often missing, so the discarded kinetic
information is not determined by the local fluid state, and the closure
becomes nonlocal, anisotropic, and tied to wave--particle resonance.

\begin{thebibliography}{99}
\bibitem{ChapmanCowling}
S.~Chapman and T.~G. Cowling,
\emph{The Mathematical Theory of Non-uniform Gases},
3rd ed., Cambridge University Press (1970).
\bibitem{Grad49}
H.~Grad,
\emph{On the kinetic theory of rarefied gases},
Commun.\ Pure Appl.\ Math.\ \textbf{2}, 331 (1949).
\bibitem{Levermore96}
C.~D. Levermore,
\emph{Moment closure hierarchies for kinetic theories},
J.\ Stat.\ Phys.\ \textbf{83}, 1021 (1996).
\bibitem{Pope}
S.~B. Pope,
\emph{Turbulent Flows},
Cambridge University Press (2000).
\bibitem{Zwanzig}
R.~Zwanzig,
\emph{Nonequilibrium Statistical Mechanics},
Oxford University Press (2001).
\bibitem{SH53}
L.~Spitzer and R.~H\"arm,
\emph{Transport phenomena in a completely ionized gas},
Phys.\ Rev.\ \textbf{89}, 977 (1953).
\bibitem{CGL56}
G.~F. Chew, M.~L. Goldberger, and F.~E. Low,
\emph{The Boltzmann equation and the one-fluid hydromagnetic equations in
the absence of particle collisions},
Proc.\ R.\ Soc.\ Lond.\ A \textbf{236}, 112 (1956).
\bibitem{Braginskii65}
S.~I. Braginskii,
\emph{Transport processes in a plasma},
in Reviews of Plasma Physics, Vol.~1, p.~205 (1965).
\bibitem{LMV83}
J.~F. Luciani, P.~Mora, and J.~Virmont,
\emph{Nonlocal heat transport due to steep temperature gradients},
Phys.\ Rev.\ Lett.\ \textbf{51}, 1664 (1983).
\bibitem{SNB00}
G.~P. Schurtz, Ph.~D. Nicola\"i, and M.~Busquet,
\emph{A nonlocal electron conduction model for multidimensional radiation
hydrodynamics codes},
Phys.\ Plasmas \textbf{7}, 4238 (2000).
\bibitem{HP90}
G.~W. Hammett and F.~W. Perkins,
\emph{Fluid moment models for Landau damping with application to the
ion-temperature-gradient instability},
Phys.\ Rev.\ Lett.\ \textbf{64}, 3019 (1990).
\bibitem{Hunana19}
P.~Hunana et al.,
\emph{An introductory guide to fluid models with anisotropic temperatures.
Part 2. Kinetic theory, Pad\'e approximants and Landau fluid closures},
\href{https://arxiv.org/abs/1901.09360}{arXiv:1901.09360} (2019).
\end{thebibliography}

\end{document}
